% ---
% title: Two borrowers, independent defaults
% tags: finance, independent-rv, exponential
% ---

\begin{problem}
	\begin{context}
		A bank's credit loss depends on how many of its borrowers default
		(stop paying their debt), not only on each borrower alone. The simplest
		model of a portfolio assumes that the default times of the borrowers are
		independent: one company's default tells us nothing about the other's.
		Each default time is modeled with an exponential distribution, whose
		rate is the yearly hazard rate of the company.
	\end{context}
	A bank has lent to two companies, A and B. Their default times $T_A$ and
	$T_B$, in years, are independent and exponential with rates $0.02$ and
	$0.03$, so their densities are
	\[ f_{T_A}(t)=0.02\,e^{-0.02t},\qquad f_{T_B}(t)=0.03\,e^{-0.03t},
		\qquad t\ge0. \]
	(a) Find the probability that A survives 10 years (does not default in
	the first 10 years), and the same for B.
	(b) What is the probability that both companies survive 10 years?
	(c) What is the probability that at least one of them defaults within
	10 years?
	(d) What is the probability that both default within 10 years?
\end{problem}

\begin{slideproblem}
	The default times $T_A$, $T_B$ (years) of two companies are
	\emph{independent}, exponential with rates $0.02$ and $0.03$.
	(a) Find $P(T_A>10)$ and $P(T_B>10)$.
	(b) Probability that both survive 10 years?
	(c) Probability that at least one defaults within 10 years?
	(d) Probability that both default within 10 years?
\end{slideproblem}

\begin{slides}
	\small
	\exstatement
	\pause
	\sol (a) Exponential with rate $\lambda$: $F(t)=1-e^{-\lambda t}$, so
	\[ P(T>t)=1-F(t)=e^{-\lambda t}. \]
	\pause
	\vspace{-2mm}
	\[ P(T_A>10)=e^{-0.02\cdot10}=e^{-0.2}\approx\alert{0.819},\qquad
		P(T_B>10)=e^{-0.03\cdot10}=e^{-0.3}\approx\alert{0.741}. \]
	\pause
	(b) $T_A$, $T_B$ independent: probabilities of events about them
	\alert{multiply}:
	\[ P(T_A>10,\ T_B>10)=P(T_A>10)\,P(T_B>10)=e^{-0.2}e^{-0.3}
		=e^{-0.5}\approx\alert{0.607}. \]
\end{slides}

\begin{slides}
	\small
	\begin{columns}[T]
		\begin{column}{.56\textwidth}
			(c) ``At least one defaults'' is the complement of ``both
			survive'':
			\[ 1-e^{-0.5}\approx1-0.607=\alert{0.393}. \]
			\uncover<2->{(d) $P(T_A\le10)=1-e^{-0.2}\approx0.181$,\\
			$P(T_B\le10)=1-e^{-0.3}\approx0.259$.}
			\uncover<3->{\[ P(T_A\le10,\ T_B\le10)\approx0.181\cdot0.259\approx\alert{0.047}. \]}
			\uncover<4->{\alert{In a crisis} defaults come together: independence
			\emph{underestimates} (d).}
		\end{column}
		\begin{column}{.40\textwidth}
			\centering
			\begin{tikzpicture}[scale=.5]
				\fill[teal!20] (2.5,2.5) rectangle (6,6);
				\uncover<3->{\fill[gold!50] (0,0) rectangle (2.5,2.5);}
				\draw[->] (0,0) -- (6.3,0) node[below left,font=\scriptsize] {$t_A$};
				\draw[->] (0,0) -- (0,6.3) node[below left,font=\scriptsize] {$t_B$};
				\draw[dashed] (2.5,0) -- (2.5,6) (0,2.5) -- (6,2.5);
				\node[below,font=\scriptsize] at (2.5,0) {$10$};
				\node[left,font=\scriptsize] at (0,2.5) {$10$};
				\node[font=\scriptsize,align=center] at (4.25,4.25) {both\\survive};
				\uncover<3->{\node[font=\scriptsize,align=center] at (1.25,1.25) {both\\def.};}
			\end{tikzpicture}
		\end{column}
	\end{columns}
\end{slides}

\begin{solution}
	\idea Each default time is exponential, so its survival probability is
	$P(T>t)=e^{-\lambda t}$. Because $T_A$ and $T_B$ are independent, the
	probability that $T_A$ falls in one set \emph{and} $T_B$ falls in
	another is the product of the two probabilities. The complement rule
	gives ``at least one''.

	\step{(a)} First the distribution function of an exponential random
	variable with rate $\lambda$ (see also Exercise~\exref{default-time}). For
	$t\ge0$ an antiderivative of $\lambda e^{-\lambda s}$ is $-e^{-\lambda s}$, so
	\[ F(t)=\int_{0}^{t}\lambda e^{-\lambda s}\,ds
		=\Bigl[-e^{-\lambda s}\Bigr]_{0}^{t}
		=-e^{-\lambda t}-\bigl(-e^{0}\bigr)=1-e^{-\lambda t}. \]
	``Survives 10 years'' means the default time is later than 10, and
	\[ P(T>t)=1-P(T\le t)=1-F(t)=e^{-\lambda t}. \]
	With $\lambda=0.02$ and $\lambda=0.03$:
	\[ P(T_A>10)=e^{-0.02\cdot10}=e^{-0.2}\approx0.819,\qquad
		P(T_B>10)=e^{-0.03\cdot10}=e^{-0.3}\approx0.741. \]
	Company B is riskier (higher rate), so it is less likely to survive.

	\step{(b)} ``Both survive'' is the event $\{T_A>10,\ T_B>10\}$: an
	event about $T_A$ \emph{and} an event about $T_B$ at the same time. The
	random variables are independent, which means exactly that for any sets
	$C$ and $D$
	\[ P(T_A\in C,\ T_B\in D)=P(T_A\in C)\,P(T_B\in D). \]
	With $C=D=(10,\infty)$:
	\[ P(T_A>10,\ T_B>10)=P(T_A>10)\,P(T_B>10)=e^{-0.2}\cdot e^{-0.3}
		=e^{-0.5}\approx0.607. \]
	(The exponents add, because $e^{x}e^{y}=e^{x+y}$.)

	\step{(c)} ``At least one defaults within 10 years'' happens exactly when
	it is \emph{not} true that both survive. So it is the complement of the
	event in (b):
	\[ P(\text{at least one defaults})=1-P(T_A>10,\ T_B>10)
		=1-e^{-0.5}\approx1-0.607=0.393. \]

	\step{(d)} The default probabilities within 10 years are the complements
	of (a):
	\[ P(T_A\le10)=1-e^{-0.2}\approx0.181,\qquad
		P(T_B\le10)=1-e^{-0.3}\approx0.259. \]
	By independence again, now with $C=D=[0,10]$,
	\[ P(T_A\le10,\ T_B\le10)=\bigl(1-e^{-0.2}\bigr)\bigl(1-e^{-0.3}\bigr)
		\approx0.1813\cdot0.2592\approx0.047. \]
	In the $(t_A,t_B)$-plane the four events are four regions cut out by the
	lines $t_A=10$ and $t_B=10$: ``both survive'' is the upper right region,
	``both default'' is the square $[0,10]\times[0,10]$, and ``at least one
	defaults'' is everything except the upper right region.
	\begin{center}
		\begin{tikzpicture}[scale=.55]
			\fill[teal!20] (2.5,2.5) rectangle (6,6);
			\fill[gold!50] (0,0) rectangle (2.5,2.5);
			\draw[->] (0,0) -- (6.3,0) node[below left,font=\small] {$t_A$};
			\draw[->] (0,0) -- (0,6.3) node[below left,font=\small] {$t_B$};
			\draw[dashed] (2.5,0) -- (2.5,6) (0,2.5) -- (6,2.5);
			\node[below,font=\small] at (2.5,0) {$10$};
			\node[left,font=\small] at (0,2.5) {$10$};
			\node[font=\scriptsize,align=center] at (4.25,4.25) {both survive\\$\approx0.607$};
			\node[font=\scriptsize,align=center] at (1.25,1.25) {both\\default\\$\approx0.047$};
		\end{tikzpicture}
	\end{center}
	\begin{answer}
		(a) $P(T_A>10)=e^{-0.2}\approx0.819$, $P(T_B>10)=e^{-0.3}\approx0.741$
		\quad (b) $e^{-0.5}\approx0.607$ \quad (c) $1-e^{-0.5}\approx0.393$
		\quad (d) $(1-e^{-0.2})(1-e^{-0.3})\approx0.047$
	\end{answer}
	\begin{remark}
		Independence is a modeling assumption, and for defaults it is often
		wrong. In a recession many companies get into trouble at the same time,
		so defaults come together. Then the probability that both default is
		larger than the product $0.047$: the independent model
		\emph{underestimates} the chance of joint defaults, which are the
		events that cause the largest losses for the bank.
	\end{remark}
\end{solution}
